\SourceLeaf{67}
$\delta$) Est igitur, si ponatur

\[
\left(\frac{1-\BareRoot{(1-nn)}}{n}\right) = a,
\]

\[
\begin{aligned}
\log(1+n\cos\varphi) & = -\log\frac{2a}{n} + \frac{2}{1}a\cos\varphi - \frac{2}{2}a^2\cos2\varphi +\\
&{} + \frac{2}{3}a^3\cos3\varphi + .\,.\,.\,. + (-1)^{h+1}\frac{2}{h}a^h\cos h\varphi + .\,.\,.\,.
\end{aligned}
\]

\textls[140]{Adnotatio.} Formula haec ad functionum trigonometricarum logarithmos obtinendos summi
usus est.

Sit enim $n=1$, fiet $a=1$, et habebitur

\[
\log(1+\cos\varphi) = -\log 2 + \frac{2}{1}\cos\varphi - \frac{2}{2}\cos2\varphi + \frac{2}{3}\cos3\varphi - \frac{2}{4}\cos4\varphi + .\,.\,.,
\]

\noindent quae series, cum sit $1+\cos\varphi = 2\cos\frac{1}{2}\varphi^2$, 
dat quoque

\[
\log\cos\tfrac{1}{2}\varphi = -\log 2 + \cos\varphi - \tfrac{1}{2}\cos2\varphi + \tfrac{1}{3}\cos3\varphi - \tfrac{1}{4}\cos4\varphi + .\,.\,.
\]

\noindent adeoque, substituendo $2\varphi$ in locum ipsius $\varphi$

\[
\log\cos\varphi = -\log 2 + \cos2\varphi - \tfrac{1}{2}\cos4\varphi + \tfrac{1}{3}\cos6\varphi - \tfrac{1}{4}\cos8\varphi + .\,.\,.\,.
\]

\noindent Similiter, si ponatur $n=-1$, ut fiat $a=-1$, habebitur

\[
\log(1-\cos\varphi) = -\log 2 - \frac{2}{1}\cos\varphi - \frac{2}{2}\cos2\varphi - \frac{2}{3}\cos3\varphi - \frac{2}{4}\cos4\varphi - .\,.\,.\,.
\]

\noindent adeoque, cum sit $1-\cos\varphi = 2\sin\tfrac{1}{2}\varphi^2$, 

\[
\log\sin\tfrac{1}{2}\varphi = -\log 2 - \cos\varphi - \tfrac{1}{2}\cos2\varphi - \tfrac{1}{3}\cos3\varphi - \tfrac{1}{4}\cos4\varphi - .\,.\,.\,.
\]

\noindent aut, ut antea, $2\varphi$ in locum ipsius $\varphi$ substituendo

\[
\log\sin\varphi = -\log 2 - \cos2\varphi - \tfrac{1}{2}\cos4\varphi - \tfrac{1}{3}\cos6\varphi - \tfrac{1}{4}\cos8\varphi - .\,.\,.\,.
\]

\noindent Sequitur ex his, cum sit $\log\operatorname{tang}\varphi = \log\sin\varphi - \log\cos\varphi$

\[
\log\operatorname{tang}\varphi = -2\left(\cos2\varphi + \tfrac{1}{3}\cos6\varphi + \tfrac{1}{5}\cos10\varphi + \tfrac{1}{7}\cos14\varphi + .\,.\,.\right)
\]

\noindent et ex eodem fonte ceterarum functionum trigonometricarum logarithmi, licet per series
non celeriter convergentes, habebuntur.

$\varepsilon$) Cum sit

\[
\log(1+n\cos\varphi) = -\log\frac{2a}{n} + \frac{2}{1}a\cos\varphi - \frac{2}{2}a^2\cos2\varphi + \frac{2}{3}a^3\cos3\varphi - .\,.\,.\,.\,.,
\]

\noindent integratio formulae $\log(1+n\cos\varphi).\,d\varphi$ in aperto est, fitque

\[
\int d\varphi\,\log(1+n\cos\varphi) = -\left(\log\frac{2a}{n}\right)\varphi + \frac{2}{1}a\sin\varphi - \frac{2}{4}a^2\sin2\varphi + \frac{2}{9}a^3\sin3\varphi - .\,.\,. 
\]

